\chapter{Conclusion}
\label{chap:conclusion}

\section{Summary of Findings}
\label{sec:concl-summary}
% Key results across all three evaluation levels
% Direct answer to the research question: how consistent are sovereign LLMs
%   in a structured therapeutic protocol -- and what does that mean clinically?

This thesis asked how consistent \gls{llm} therapeutic responses are
across independent runs within a structured \gls{irt} protocol. Three
evaluation methods were applied to ten models across 6{,}000 trials
spanning five temperature conditions.

The main findings are:

\begin{enumerate}
  \item \textbf{Consistency varies widely across models and temperatures.}
    Mean Jaccard ranges from $J = 0.56$ (\modelname{Qwen~3.5~27B} at
    $T = 0.6$) to $J = 1.0$ (\modelname{OLMo~3.1~32B} at $T = 0.0$).
    Temperature is the primary driver of variation for open-weight models.

  \item \textbf{The \gls{eu}-sovereign flagship is close to proprietary
    models at recommended temperatures.} \modelname{Mistral Large~3}
    achieves $J = 0.848$ at $T = 0.15$, compared to $J = 0.871$ for
    \modelname{GPT-5.4} and $J = 0.907$ for \modelname{Sonnet~4.6}. The
    gap widens at higher temperatures.

  \item \textbf{Instability sits in strategy selection, not execution.}
    Temperature reduces plan consistency but not plan-response coherence.
    Models pick different strategies at higher temperatures but carry out
    whichever strategy they pick equally well.

  \item \textbf{Greedy decoding is not always optimal.} Three models achieve
    higher consistency at $T = 0.075$ than at $T = 0.0$. Whether this
    reflects a genuine benefit of low sampling variance or an artefact of
    the 20-trial sample size cannot be determined from these results alone.

  \item \textbf{No single metric captures the full picture.}
    The three evaluation levels are only moderately correlated
    ($\rho = 0.217$ to $0.508$). BERTScore misses strategy changes.

  \item \textbf{Patient profile may matter.} The vignette effect on
    Jaccard is not significant ($p = 0.094$), but BERTScore shows a
    marginal effect ($p = 0.033$). Per-vignette reporting can reveal
    weaknesses that aggregate scores mask.
\end{enumerate}

Regarding the two contributions stated in \autoref{sec:intro-objectives}:
the three-level framework (Contribution~1) fulfilled its design goal by
identifying that instability is concentrated at the strategy selection
layer, a finding that neither BERTScore nor coherence alone could have
revealed. The cross-model comparison (Contribution~2) shows that
\modelname{Mistral Large~3} approaches proprietary-level stability at its
recommended operating temperature, though a small gap remains.


\section{Limitations}
\label{sec:concl-limitations}
% Scope: single protocol, single phase (IRT rescripting only)
% Simulation: synthetic patients differ from real clinical interaction in
%   ways that cannot be fully controlled
% Six vignettes sample the difficulty range but cannot claim exhaustive
%   coverage of patient variability
% Method 3 is exploratory -- judge reliability is unverified and should
%   not be treated as ground truth
% Fixed taxonomy constrains the measurable strategy space -- unmeasured
%   strategies may show different stability patterns
% Computational stability does not equal clinical safety

The study evaluates a single phase (rescripting) of a single protocol
(\gls{irt}). The rescripting phase was selected because it requires
active strategy decisions, but stability in other stages or other
therapeutic protocols may differ.

All patient interactions are simulated. Synthetic patients follow scripted
behavioural profiles and do not exhibit the unpredictable responses of
real patients. Six profiles sample the difficulty range
(\autoref{sec:meth-dialogue}) but cannot claim exhaustive coverage. No
human participants were involved in this study, and the results measure
computational consistency rather than clinical safety or efficacy.
\gls{irt} targets nightmares rather than psychiatric diagnoses, which
limits clinical risk, but any deployment to real patients would require
validation beyond what this framework provides.

The \gls{llm} judge (\modelname{Gemini 3.1 Pro}) was validated on 60
stratified trials with substantial human agreement ($\kappa = 0.777$),
but this covers only 1\% of the 6{,}000 total trials. A larger
annotation study would strengthen confidence in the coherence scores.

The six-category strategy taxonomy constrains what can be measured.
Strategies that fall outside these categories are invisible to Method~1,
and different category boundaries would produce different absolute
consistency scores.

Finally, consistency is not the same as quality. A model that consistently
selects an inappropriate strategy scores high on plan consistency.
Stability is necessary for clinical deployment but not sufficient on its
own.


\section{Future Work}
\label{sec:concl-future}
% Full IRT protocol evaluation -- beyond the rescripting phase
% Extension to other manualized therapies (CBT, DBT, PE)
% Longitudinal drift: does stability degrade across model version updates?
% Multi-language evaluation -- German therapeutic context for reDreamAI
% Clinical validation: human therapist ratings as external ground truth
% Replication with future model generations to track whether stability
%   improves across release cycles

Several directions follow from this work.

Extending the evaluation beyond the rescripting phase to all \gls{irt}
stages would test whether the observed patterns hold when the
model faces different types of clinical decisions. More broadly, the
three-level framework applies to other therapy protocols where the
therapist chooses from a defined set of techniques, such as \gls{cbt} or
motivational interviewing. Each would need its own strategy taxonomy, but
the evaluation setup carries over.

Since model vendors update weights and infrastructure without notice,
repeating the evaluation across model versions would show whether
stability improves or fluctuates over time. reDreamAI targets
German-speaking users, so evaluating stability in German would test
whether the English-language patterns transfer.

Finally, newer model architectures may affect output consistency
independently of temperature. Whether architectural choices such as
hybrid attention mechanisms can improve therapeutic consistency beyond
what temperature tuning achieves remains an open question.
